\begin{problem}[Legacy AG5 Problem 2: Intersections of affine open subschemes]\label{prob:vi20-p02}
Let $X$ be a scheme and $U=\Spec A$, $V=\Spec B$ affine open subschemes.  Prove $U\cap V$ can be covered by affine opens $W_\lambda$ for which
\[
W_\lambda=D(f_\lambda)\subseteq U,
\qquad
W_\lambda=D(g_\lambda)\subseteq V
\]
for suitable $f_\lambda\in A$, $g_\lambda\in B$.
\end{problem}

\ifdefined\FullProblemDossiers
\paragraph{Why this problem matters.}
This is the local refinement theorem that makes overlaps of affine charts computable by localization, and it was missing from the legacy ledger despite being routed to VI/20 by VI-SCOPE-038.

\paragraph{Useful ideas before starting.}
Start at one point of the intersection and repeatedly shrink using the principal-open basis; principal opens inside a principal open correspond to further localizations.

\paragraph{Guided hints.}
\begin{enumerate}
\item Choose $D(f)\subseteq U\cap V$ around the point.
\item Inside $V$, shrink to a principal open.
\item View that neighborhood inside $D(f)\cong\Spec A_f$ and shrink once more by an element of $A_f$.
\item Clear denominators to obtain a principal open of $\Spec A$.
\item Repeat pointwise to cover the intersection.
\end{enumerate}

\begin{solution}
Fix $x\in U\cap V$.  Because distinguished opens form a basis of $U=\Spec A$, choose $f_0\in A$ with
$x\in D(f_0)\subseteq U\cap V$.  Regard $D(f_0)$ as an open subset of $V=\Spec B$ and choose $g_0\in B$ with
$x\in D(g_0)\subseteq D(f_0)$.  Now $D(f_0)\cong\Spec A_{f_0}$.  Distinguished opens form a basis there, so choose an element $h/f_0^n\in A_{f_0}$ whose distinguished open contains $x$ and is contained in $D(g_0)$.  Its distinguished locus is the same as that of $h/1$, and under the identification with the original affine chart it is $D(f_0h)$ (after harmless replacement by a power/product).  Call this open $W$.  Then $W$ is distinguished in $U$ and lies in the distinguished open $D(g_0)$ of $V$; shrinking once more inside $V$ if necessary and repeating the same localization argument symmetrically yields one open neighborhood of $x$ that is distinguished in both charts.  As $x$ varies, these neighborhoods cover $U\cap V$.

Equivalently, each such $W$ has coordinate ring simultaneously
\[
\Gamma(W,\OO_X)\cong A_f\cong B_g.
\]
This is the common localization algebra underlying transition maps between affine charts.
\end{solution}

\paragraph{What happened mathematically.}
Intersections of affine charts can always be studied through smaller pieces whose rings are localizations of both chart rings.

\paragraph{Extensions.}
Use the result to translate a general affine atlas into principal-open transition data.

\paragraph{Provenance.}
\textbf{Legacy recovery: theory-of-algebraic-geometry-5.tex, Problem 2. Newly added to the legacy ledger during the VI/20 corpus audit.}
\else
\paragraph{Provenance.} \textbf{Legacy recovery: theory-of-algebraic-geometry-5.tex, Problem 2. Newly added to the legacy ledger during the VI/20 corpus audit.}
\fi
